% The post src/highlights/diy-mea.md as a paper. The prose is the post's; only the
% typesetting (units, figures, the reagent table, references) is new.
% Build with ./build.sh, which also fills .cache/img with jpg copies of the images.
\documentclass[11pt,letterpaper]{article}

\usepackage[margin=1.15in,top=1.1in,bottom=1.2in]{geometry}
\usepackage{fontspec}
\usepackage{microtype}
\usepackage{graphicx}
\usepackage{float} % [H] keeps each figure beside the paragraph that talks about it
\usepackage{xcolor}
\usepackage{siunitx}
\usepackage{booktabs}
\usepackage{tabularx}
\usepackage{enumitem}
\usepackage[font=small,labelfont=bf,skip=6pt]{caption}
\usepackage{subcaption}
\usepackage{titlesec}
\usepackage{fancyhdr}
\usepackage[colorlinks=true,allcolors=gardenviolet]{hyperref}

\definecolor{gardenviolet}{HTML}{6D28D9}
\graphicspath{{.cache/img/}}
\newcommand{\photoheight}{5.4cm}
\sisetup{group-separator={\,},group-minimum-digits=5}
\setlist{itemsep=2pt,topsep=4pt}
\setlength{\parskip}{0.55em}
\setlength{\parindent}{0pt}
\titleformat{\section}{\large\bfseries}{\thesection}{0.7em}{}
\titleformat{\subsection}[runin]{\bfseries}{}{0pt}{}[.]
\titlespacing*{\subsection}{0pt}{0.9em}{0.5em}

\pagestyle{fancy}
\fancyhf{}
\renewcommand{\headrulewidth}{0pt}
\fancyfoot[L]{\footnotesize\color{gray}\href{https://adiabatic.garden/highlights/diy-mea/}{adiabatic.garden/highlights/diy-mea}}
\fancyfoot[R]{\footnotesize\color{gray}\thepage}

\hypersetup{
  pdftitle={Making low-cost microelectrode arrays},
  pdfauthor={noesis and Eigenlucy},
}

\begin{document}

\begin{center}
  {\LARGE\bfseries Making low-cost microelectrode arrays\par}
  \vspace{0.9em}
  {\large noesis and Eigenlucy\par}
  \vspace{0.3em}
  {\small\color{gray} 2025 \quad\textperiodcentered\quad lab equipment, neuroscience, experimentalist-confetti\par}
\end{center}

\vspace{0.4em}
\begin{quote}
\small
A single microelectrode array of 59 channels costs \$300, and integrated neuronal
headstage setups cost \$10k--\$50k. We experimented with building a cheap MEA using
three laser deposition techniques: laser-induced graphene, aluminum vapour deposition
and copper metalization. So far the graphene version costs $<$\$1 and \SI{5.5}{\minute}
per unit.
\end{quote}
\vspace{0.2em}

\section{Motivation: Neuronal headstages bottlenecked researchers}

2025 was full of startup hype about neuron computers. At Frontier Tower's neurotech
floor, friends were discussing biological neuronal responses to nootropics and
neuropeptides. However, the neuronal headstages were hard to get. Student clubs e.g.\
UC Berkeley didn't rent them out, nor did neurotech startups e.g.\ Science.xyz.

For reference, a single microelectrode array of 59 channels from multichannel systems
costs \$300.

Not to mention, integrated neuronal headstage setups from Axion biosystems cost
\$10k--\$50k.

Thus, we experimented with building a cheap MEA, and so far graphene version costs
$<$\$1 and \SI{5.5}{\minute} per unit.

\section{We tested 3 laser deposition techniques}

\begin{enumerate}
  \item Laser-induced graphene
  \item Aluminum vapour deposition
  \item Copper metalization
\end{enumerate}

We've looked into laser ablation: Indium tin oxide. This was not tested but may also
be worth a shot.

\section{Designing an MEA}

\textbf{Iteration 1: OnShape.} Created at $\sim$12 am from circles and linear patterns.
I found out this shortcircuited, since I naively traced the well as a ring, so I
created version 2!

\textbf{Iteration 2.} The improved MEA in KiCad. This dotted grid directly come into
contact with neurons, and extend out to the peripherals.

\textbf{Iteration 3.} Using existing MEA design files provided by
Mind-in-Vitro~\cite{miv}, openable in Klayout.

\begin{figure}[H]
  \centering
  \begin{subfigure}[b]{0.32\textwidth}
    \includegraphics[width=\linewidth,height=3.6cm,keepaspectratio]{onshape}
    \caption{Iteration 1: OnShape}
  \end{subfigure}\hfill
  \begin{subfigure}[b]{0.32\textwidth}
    \includegraphics[width=\linewidth,height=3.6cm,keepaspectratio]{mea-routes}
    \caption{Iteration 2: KiCad}
  \end{subfigure}\hfill
  \begin{subfigure}[b]{0.32\textwidth}
    \includegraphics[width=\linewidth,height=3.6cm,keepaspectratio]{klayout}
    \caption{Iteration 3: Klayout}
  \end{subfigure}
  \caption{Three iterations of the MEA design.}
  \label{fig:design}
\end{figure}

\section{Most promising: Laser-induced graphene}

At atopile's machine shop, we directed a laser beam at the polyimide tape, for graphene
to become deposited on the surface (Figure~\ref{fig:lig}). We are using microscope
slides from frontier tower's biopunk lab run by Elliot Roth. Techniques were adapted
from this porous graphene paper~\cite{lig}.

The laser's energy causes lattice vibrations within the polyimide: Under high surface
temperatures of $>$\SI{2500}{\celsius}, C=O, C--N, and C--O bonds are broken, freeing
non-carbon gases from this tape. The carbon atoms rearrange into graphene with its
hexagonal structure.

We tuned a CO\textsubscript{2} laser operating at \SI{10600}{\nano\metre} to
\SI{5.5}{\watt} with a scanning rate of \SI{88}{\milli\metre\per\second} with
100~lines/cm and directed the beam at the kapton tape.

The centre-to-centre distance between electrodes: \SI{0.15}{\centi\metre}, with a
\SI{0.25}{\centi\metre} electrode radius. The trace width: \SI{0.2}{\milli\metre}.

A hurdle is creating continuously conducting traces.

\begin{figure}[H]
  \centering
  \begin{subfigure}[b]{0.3\textwidth}
    \centering\includegraphics[height=\photoheight]{beautiful-lig}
    \caption{Graphene deposited on polyimide tape}
  \end{subfigure}\hspace{0.04\textwidth}
  \begin{subfigure}[b]{0.3\textwidth}
    \centering\includegraphics[height=\photoheight]{lig}
    \caption{The triple-layered electrode version}
  \end{subfigure}
  \caption{Laser-induced graphene.}
  \label{fig:lig}
\end{figure}

\subsection{Trace width microadjustments}
Given that SVGs are a set of lines defined using the XML format, we just made them
thicker. However, in practice, laser software, e.g.\ Xtool or Lightburn, treat all
lines as infinitely thin.

We then tried blowing up the JPGs, which are bitmaps; However, a change from
\SI{0.2}{\milli\metre} to \SI{0.3}{\milli\metre} in the JPG led to a fuzzy visual
appearance and shortcircuited. A lazy hack we found was just to duplicating the traces
three times in the laser software from Xtool. Additionally, we varied the scan rate
($\pm$\SI{20}{\milli\metre\per\second}) and lines per cm (max at 200, min at 50) and set
the laser scanning to be monodirectional.

\subsection{Examining raw trail conductivity}
The triple-layered electrode version was consistently conductive. However, the graphene
easily wears off, or encounters insulation at corners.

The next step would be more functional testing (saline) and compare to the standard
Utah array.

\section{Accidental discovery: Aluminum vapour deposition}

While implementing laser-induced graphene, we noticed dense metallic stripes appeared
on glass---This certainly didn't look like graphene (Figure~\ref{fig:aluminum}). So we
peeled off the kapton tape, put the glass slide into the laser chamber and turned the
laser on again. This was in fact, the vapour from the underlying sheet of aluminum
which had been sitting inside the glass in a diode-pumped solid state laser (DPSS) of
\SI{1064}{\nano\metre}
(\href{https://en.wikipedia.org/wiki/Physical_vapor_deposition}{physical vapour deposition}).

\begin{figure}[H]
  \centering
  \includegraphics[height=\photoheight]{aluminum}
  \caption{Aluminum deposited on a glass slide.}
  \label{fig:aluminum}
\end{figure}

\section{Error-prone: Copper metalization}

Review: In the traditional electroplating process, an object, the cathode, becomes
coated with a thin layer of metal (Figure~\ref{fig:plating}).

The cathode and the anode are initially submerged in ionic solution and a circuit is
completed. When the battery applies its gradient, two processes happen:
\begin{enumerate}
  \item Chemical. The +ve ions start migrating to the cathode to coat it
  \item Electrical. The cathode ``craves'' electrons and the electrons move to the
        cathode. Simplified copper reaction:
        $\mathrm{Cu^{2+} + 2\,e^{-} \rightarrow Cu(s)}$.
\end{enumerate}

In copper plating, a focused laser beam mimicks the voltage differential like a
battery. As it is directed at the copper solution, it causes a thermal gradient: Heated
and unheated surfaces. The heated area gains copper under the plate.

\begin{figure}[H]
  \centering
  \begin{subfigure}[b]{0.3\textwidth}
    \centering\includegraphics[height=\photoheight]{copper-plating}
    \caption{Traditional electroplating}
    \label{fig:plating}
  \end{subfigure}\hspace{0.04\textwidth}
  \begin{subfigure}[b]{0.3\textwidth}
    \centering\includegraphics[height=\photoheight]{copper-completed}
    \caption{Results of the deposition}
    \label{fig:copper-result}
  \end{subfigure}
  \caption{Copper metalization.}
\end{figure}

\subsection{Preparing the electroplating solution}
Add water first, and then add chemicals slowly from least reactive to most. i.e.\ Add
distilled water, followed by alcohol, glycerin, NaOH and then CuSO\textsubscript{4}
(Table~\ref{tab:solution}).

\begin{table}[H]
  \centering
  \small
  \newcolumntype{R}[1]{>{\raggedright\arraybackslash}p{#1}}
  \begin{tabularx}{\textwidth}{@{}l R{2.5cm} R{1.9cm} R{2.4cm} >{\raggedright\arraybackslash}X@{}}
    \toprule
    Molarity & Grams/mol & Grams (to solve for) & Chemical & Description \\
    \midrule
    \SI{0.4}{M}   & \SI{92.1}{g/mol}    & \SI{3.68}{g}  & Glycerin & Used as a substitute to Rochelle's salt \\
    \SI{0.125}{M} & \SI{39.997}{g/mol}  & \SI{0.5}{g}   & NaOH & Makes it an alkaline solution \\
    \SI{0.1}{M}   & \SI{159.609}{g/mol} & \SI{1.6}{g}   & CuSO\textsubscript{4} & Copper sulphate: copper ion donor \\
    \SI{6}{M}     & \SI{46.086}{g/mol} (pure ethanol) & \SI{27.65}{g} & C\textsubscript{2}H\textsubscript{5}OH (Ethanol) & Acts as a reducing agent. In practice, ethanol is often diluted to 70\% \\
    \SI{55.5}{M}  & \SI{18.02}{g/mol}   & \SI{100}{g}   & Distilled water & To create the solution \\
    \bottomrule
  \end{tabularx}
  \caption{The electroplating solution.}
  \label{tab:solution}
\end{table}

We placed the solvent on a hot plate, with a magnetic rod for stirring
(Figure~\ref{fig:bench}).

\begin{figure}[H]
  \centering
  \begin{subfigure}[b]{0.3\textwidth}
    \centering\includegraphics[height=\photoheight]{copper-waterbath}
    \caption{The bench}
  \end{subfigure}\hspace{0.04\textwidth}
  \begin{subfigure}[b]{0.3\textwidth}
    \centering\includegraphics[height=\photoheight]{copper-heating}
    \caption{On the hot plate}
  \end{subfigure}
  \caption{Preparing the electroplating solution.}
  \label{fig:bench}
\end{figure}

The biggest challenge here was when the scale started tweaking. We became suspicious
when the \SI{100}{\milli\litre} of water weighed 125 grams. Consequentially, in our 2nd
copper replication attempt, the copper solution instantly precipitated into copper
hydroxide---a cloudy, turquoise colour, not what we want! The ideal solution should
look transparent, dark blue.

We tried a series of quick troubleshooting e.g.\ switching power supplies, using
different scales, but there weren't matching power supplies of \SI{9}{\volt} for any
other scale in the lab.

Thus, I poured in some ethanol (\SI{1}{\milli\litre} of ethanol was 0.789 grams) to see
how much the scale was off by---approximately 29.5875 grams. From then on, we performed
everything according to the scale's offset mass (1.25x). This method has been
time-intensive and prone to error.

\section{Not tested: PEDOT:PSS trails}

Motivation. Researchers may use red and green fluorescent proteins in their living
neurons which includes filtering colours. So the yellow polyimide tape may present dark
under a microscope. They have a perfect reason to want transparent MEAs?

Here are some materials to use (loosely inspired by this paper)
\begin{itemize}
  \item Indium tin oxide coated borosilicate glass as substrate
  \item Brush-on negative insulator (1002F epoxy) -- To cover the non-conductive regions
  \item PEDOT:PSS (PH1000 + 5\,\% DMSO or ethylene-glycol) -- Making the surface
        biocompatible, although ITO is already safe.
\end{itemize}

Then, the plasma engine could increase the surface energy of glass.

\section{Next steps}

It would be really good to replicate the entire neuronal headstage from Mind-in-Vitro
to broaden the accessibility. Please get in touch
\href{mailto:corpus.logica@gmail.com}{corpus.logica@gmail.com} if you want to support
this project such as by offering ephys lab spaces.

\section*{Credits}

Eigenlucy (main collaborator)

Supporters: Morgan Hough, Paul Han, Elliot Roth, Lucas Mair, James Mavo at frontier
tower, San Francisco as well as Atopile for letting us use their office.

Funding: 1517 Medici project, merge grant, Lucas Chu and individual donors.

\begin{thebibliography}{9}
\small
\bibitem{miv} X.~Zhang, Z.~Dou, S.~H.~Kim et al.
  \newblock Mind In Vitro Platforms: Versatile, Scalable, Robust, and Open Solutions to
  Interfacing with Living Neurons.
  \newblock \emph{Advanced Science} 11(11), 2306826, 2024.
  \newblock \href{https://doi.org/10.1002/advs.202306826}{doi:10.1002/advs.202306826}
\bibitem{lig} J.~Lin, Z.~Peng, Y.~Liu et al.
  \newblock Laser-induced porous graphene films from commercial polymers.
  \newblock \emph{Nature Communications} 5, 5714, 2014.
  \newblock \href{https://doi.org/10.1038/ncomms6714}{doi:10.1038/ncomms6714}
\end{thebibliography}

\end{document}
